\section{General formulation and fixed-architecture optimum}
\label{app:technical-projection}
\label{subsec:projection}

This appendix preserves the broader fixed-epoch quadratic formulation.
Its target may be non-affine, unlike the temporal composition theorem.

On a probability space $(\Omega,\mathcal F,\mathbb P)$, let
$\theta=\{\theta_t\}_{t\in\mathbb T}$ be an environment process,
with discrete or continuous time. Write
\begin{equation}
\theta_t=(\theta_{1,t},\ldots,\theta_{n,t})\in\mathbb R^m,
\quad \theta_{i,t}\in\mathbb R^{m_i},\quad \sum_i m_i=m,
\quad \mathbb E\|\theta_t\|^2<\infty.
\label{eq:state_partition}
\end{equation}
Write $\bar\theta=\mathbb E[\theta_t]$ and
$\Sigma_S=\operatorname{Cov}(\theta_t)$ for its mean and covariance at the fixed epoch.
The local state $\theta_{i,t}$ may describe channel conditions, demand,
workload, or sensing quality. This fixed-epoch formulation requires no
stationarity, Gaussianity, independence, or Markov assumptions. Agent $i$ chooses
$x_{i,t}\in\mathbb R^{d_i}$; the joint action is
$x_t=(x_{1,t},\ldots,x_{n,t})\in\mathbb R^d$, with $d=\sum_i d_i$.

The agents share a common per-stage cost
\begin{equation}
    f(x,\theta)
    =
    \frac{1}{2}x^\top Qx
    -
    b(\theta)^\top x
    +
    c(\theta),
    \label{eq:quadratic_cost}
\end{equation}
where $Q\in\mathbb{R}^{d\times d}$ is symmetric and positive definite,
$b:\mathbb{R}^{m}\rightarrow\mathbb{R}^{d}$ is measurable, and
$c:\mathbb{R}^{m}\rightarrow\mathbb{R}$ is an additive state-dependent term.
We assume $\mathbb E\|b(\theta_t)\|^2<\infty$ and
$\mathbb E|c(\theta_t)|<\infty$. The off-diagonal blocks of $Q$ couple agents'
actions, while $b(\theta)$ determines which states affect the preferred action.
This decision coupling need not coincide with proximity in $G$.

If the complete current state $\theta_t$ is available without information
constraints, the unique optimal joint action is
\begin{equation}
    x_t^\star
    \triangleq
    x^\star(\theta_t)
    =
    Q^{-1}b(\theta_t).
    \label{eq:clairvoyant}
\end{equation}
We call $x_t^\star$ the \emph{clairvoyant} or \emph{full-information}
decision.

Completing the square in \eqref{eq:quadratic_cost} gives
\begin{equation}
    f(x,\theta)-f(x^\star(\theta),\theta)
    =
    \frac{1}{2}
    \left(x-x^\star(\theta)\right)^\top
    Q
    \left(x-x^\star(\theta)\right).
    \label{eq:complete_square}
\end{equation}
Hence, for the quadratic model, excess decision cost is exactly a weighted
squared error in reproducing the full-information action.

For the affine specializations used to obtain closed forms, we write
\begin{equation}
    b(\theta)=B\theta+b_0,
    \qquad
    B\in\mathbb{R}^{d\times m},
    \qquad
    b_0\in\mathbb{R}^{d},
    \qquad
    K\triangleq Q^{-1}B\in\mathbb{R}^{d\times m}.
    \label{eq:affine_map_and_sensitivity}
\end{equation}
The matrix $B$ maps state into the linear term of the objective, while $K$ is
the resulting decision-sensitivity operator. The general projection results do
not require this affine form.



\paragraph{Centering the affine model.}
If $b(\theta)=B\theta+b_0$, write
$\theta^c=\theta-\mathbb E[\theta]$ and
$x^c=x-\mathbb E[x^\star]$. Then the centered desired action is
$x^{\star,c}=K\theta^c$, with $K$ as in
Eq.~\eqref{eq:affine_map_and_sensitivity}. Translating each observation and
each action by its known mean gives a bijection of the admissible rules
and leaves $x-x^\star$ unchanged. Thus Eq.~\eqref{eq:complete_square}
is precisely the tracking loss of the main text in centered coordinates.
For a non-affine $b$, the projection result below still holds, but the scalar
state prediction assumption need not give scalar prediction of $x^\star$.

\paragraph{Information and admissible rules.}
Let $\mathcal G_i\subseteq\mathcal F$ be the sigma-algebra generated by
agent $i$'s observations; for example, $\mathcal G_i=\sigma(Z_i)$.
Measurability here means that $u_i$ can be computed from those observations.
An architecture and its policy space are
\begin{equation}
\mathcal A=(\mathcal G_1,\ldots,\mathcal G_n),\qquad
\mathcal S(\mathcal A)=\{u:\mathbb E\|u\|^2<\infty,
\ u_i\text{ is }\mathcal G_i\text{-measurable for every }i\}.
\label{eq:policy_subspace}
\end{equation}
In the local example $\mathcal G^{\rm loc}_{i,t}=\sigma(\theta_{i,t})$.
If observations are enlarged without discarding old ones, then
\begin{equation}
\mathcal G_i\subseteq\mathcal G'_i\ (\text{all }i)
\quad\Longrightarrow\quad
\mathcal S(\mathcal A)\subseteq\mathcal S(\mathcal A'),\quad
R(\mathcal A')\le R(\mathcal A).
\label{eq:info_inclusion}
\end{equation}
This formalizes the information-inclusion argument in the body.

Let
\begin{equation}
    \mathcal{H}
    =
    L^2(
        \Omega,\mathcal{F},\mathbb{P};
        \mathbb{R}^{d}
    )
    \label{eq:hilbert_space}
\end{equation}
and equip $\mathcal{H}$ with the weighted inner product
\begin{equation}
    \langle u,v\rangle_Q
    \triangleq
    \mathbb{E}
    \left[
        u^\top Qv
    \right],
    \qquad
    \|u\|_Q^2
    \triangleq
    \langle u,u\rangle_Q.
    \label{eq:q_inner_product}
\end{equation}
Since $Q\succ0$, the induced norm is equivalent to the ordinary $L^2$ norm.

\begin{lemma}[Fixed-architecture quadratic projection; Radner-type specialization]
\label{lem:radner_projection}
For any static information architecture $\mathcal{A}$, the admissible policy set
$\mathcal{S}(\mathcal{A})$ defined in \eqref{eq:policy_subspace} is a closed
linear subspace of $\mathcal{H}$. Under the quadratic cost
\eqref{eq:quadratic_cost}, the unique optimal policy, up to almost-sure equality, implementable under
$\mathcal{A}$ is
\begin{equation}
    u_{\mathcal{A}}^\star
    =
    \Pi_{\mathcal{S}(\mathcal{A})}
    x_t^\star,
    \label{eq:optimal_projection}
\end{equation}
where $\Pi_{\mathcal{S}(\mathcal{A})}$ denotes orthogonal projection with
respect to \eqref{eq:q_inner_product}. Consequently,
\begin{equation}
    \boxed{
    R(\mathcal{A})
    =
    \frac{1}{2}
    \left\|
        x_t^\star
        -
        \Pi_{\mathcal{S}(\mathcal{A})}
        x_t^\star
    \right\|_Q^2 .
    }
    \label{eq:projection_regret}
\end{equation}
This identity specializes the classical static-team formulation
\cite{radner1962team} to the quadratic model used here.
\end{lemma}

\paragraph{Proof of Lemma~\ref{lem:radner_projection}.}
\begin{proof}
Linearity of $\mathcal{S}(\mathcal{A})$ follows because
$\mathcal{G}_i$-measurability is preserved under linear combinations.
Closedness follows because an $L^2$ limit of
$\mathcal{G}_i$-measurable random variables admits a
$\mathcal{G}_i$-measurable version.

For any admissible policy $u$, \eqref{eq:complete_square} gives
\begin{align}
    \mathbb{E}
    \left[
        f(u,\theta_t)
        -
        f(x_t^\star,\theta_t)
    \right]
    &=
    \frac{1}{2}
    \mathbb{E}
    \left[
        (u-x_t^\star)^\top
        Q
        (u-x_t^\star)
    \right]
    \nonumber\\
    &=
    \frac{1}{2}
    \|u-x_t^\star\|_Q^2.
    \label{eq:projection_proof}
\end{align}
Minimizing this expression over the closed subspace
$\mathcal{S}(\mathcal{A})$ is the standard Hilbert-space projection problem,
yielding \eqref{eq:optimal_projection} and \eqref{eq:projection_regret}.
\end{proof}

The projection error is $Q$-orthogonal to every admissible policy. In particular,
$2R(\mathcal A)=\|x_t^\star\|_Q^2-\|u_{\mathcal A}^\star\|_Q^2$,
which is often useful for computing regret.

\begin{corollary}[Coupled conditional normal equations]
\label{cor:conditional_normal_equations}
The unique team optimum in Lemma~\ref{lem:radner_projection} satisfies
\begin{equation}
    \boxed{
    \mathbb E\!\left[
      \bigl(Q u_{\mathcal A}^{\star}-b(\theta_t)\bigr)_i
      \;\middle|\;
      \mathcal G_i
    \right]
    =0,
    \qquad i=1,\ldots,n.
    }
    \label{eq:conditional_normal_equations}
\end{equation}
With $Q$ partitioned into action blocks, this is equivalently
\begin{equation}
Q_{ii}u_i^\star
+
\sum_{j\ne i}Q_{ij}\,
\mathbb E[u_j^\star\mid\mathcal G_i]
=
\mathbb E[b_i(\theta_t)\mid\mathcal G_i].
\label{eq:block_conditional_normal_equations}
\end{equation}
These are the coupled person-by-person stationarity equations associated with
the classical static-team formulation \cite{radner1962team}. Strict convexity
makes them sufficient for the unique team optimum here. With heterogeneous
information and
off-diagonal $Q$, one generally does not have
$u_i^\star=\mathbb E[x_i^\star\mid\mathcal G_i]$. When all components share a
common information sigma-algebra, deterministic $Q$ makes the joint solution
the ordinary vector conditional expectation.
\end{corollary}


\begin{proof}[Justification of the conditional normal equations]
For an admissible perturbation $h$, differentiating expected quadratic cost
at the optimum gives $\mathbb E[h^\top(Qu^\star-b(\theta_t))]=0$.
Choose perturbations supported on one agent's action block, with that block
an arbitrary square-integrable function of its observations. Conditional
expectation then gives Eq.~\eqref{eq:conditional_normal_equations}.
Conversely, these conditional equations make the linear term vanish in
\[
\mathbb E[f(u^\star+h,\theta_t)-f(u^\star,\theta_t)]
=\mathbb E[h^\top(Qu^\star-b(\theta_t))]+\tfrac12\mathbb E[h^\top Qh].
\]
Positive definiteness makes the remaining term positive for any nonzero
$L^2$ perturbation, proving sufficiency and uniqueness up to almost-sure equality.
\end{proof}

\paragraph{Common-information optimality.}
If every agent observes the same $Z$, let $m=\mathbb E[x^\star\mid Z]$.
For every admissible perturbation $h$, conditioning gives
$\mathbb E[h^\top Q(x^\star-m)]=0$. The error expansion then proves that
$m$ is the unique optimum, regardless of the off-diagonal entries of $Q$.
For different observations this argument does not apply componentwise.
\paragraph{Delayed common snapshots.}
A common old snapshot gives regret
\[
\tfrac12\mathbb E[(x_t^\star-\mathbb E[x_t^\star\mid\theta_{t-\tau}])^\top
Q(x_t^\star-\mathbb E[x_t^\star\mid\theta_{t-\tau}])]
\]
even for a non-affine target. If its conditional mean converges in mean square to the unconditional
mean, this regret converges to the no-information value. Under a stationary
Markov environment, data processing makes this prediction loss nondecreasing
in delay; arbitrary snapshot information sets themselves need not be nested.
\paragraph{Degenerate normalization.}
The best constant is $\mathbb E[x^\star]$; its regret is half the weighted
variance of $x^\star$. In centered coordinates, $R_\infty=0$ and $Q\succ0$
imply $x_t^\star=0$ almost surely. The zero rule has zero regret under
every architecture, so the unnormalized composition law remains valid
while normalized shares are undefined.

Stationarity makes expected snapshot regret time-invariant. It is not needed
for the fixed-epoch projection result. Identifying long-run sample-path
averages with these expectations additionally requires ergodicity.

\subsection{Temporal special cases}
\label{app:temporal-special-cases}
The stationary Gaussian Ornstein--Uhlenbeck process is a sufficient example
of the main conditional-mean assumption, rather than its definition:
\begin{equation}
d\theta_t=-\frac1T(\theta_t-\bar\theta)\,dt
+\sqrt{2/T}\,\Sigma_S^{1/2}dW_t.
\label{eq:OU_process}
\end{equation}
Here $W_t$ is standard multivariate Brownian motion and $\Sigma_S\succeq0$.
Its centered transition has independent Gaussian innovations, and
\begin{equation}
\operatorname{Cov}(\theta_t,\theta_s)=e^{-|t-s|/T}\Sigma_S.
\label{eq:separable_space_time_covariance}
\end{equation}
For a stationary AR(1) process with coefficient $0<c<1$ and independent,
mean-zero square-integrable innovations, the lag-$k$ prediction coefficient
is $c^k$ and the composition weights are $1-c^{2k}$ and $c^{2k}$.
Neither statement asserts that a latest neighborhood snapshot is sufficient
relative to its entire observation history.

\section{Aggregate-coupling calculations}
\label{app:technical-aggregate}
This section supplies the matrix arithmetic for Proposition~\ref{prop:aggregate_local}
above. The individual regrets are
\begin{align}
    R_\infty
    &=
    \frac{
        n\sigma^2
        \left(
            q+\gamma(n-1)
        \right)
    }{
        2q(q+\gamma n)
    },
    \label{eq:aggregate_Rinf}
    \\
    R_{\mathrm{loc}}
    &=
    \frac{
        n\sigma^2\gamma^2(n-1)
    }{
        2q(q+\gamma)(q+\gamma n)
    }.
    \label{eq:aggregate_Rloc}
\end{align}
For $Q=qI+\gamma\mathbf1\mathbf1^\top$,
\[
Q^{-1}=\frac1q\left(I-\frac{\gamma}{q+\gamma n}\mathbf1\mathbf1^\top\right).
\]
\begin{proof}
Recall the coupled normal equations \eqref{eq:block_conditional_normal_equations}
from Appendix~\ref{app:technical-projection}:
\[
 Q_{ii}u_i+\sum_{j\ne i}Q_{ij}\mathbb E[u_j\mid\mathcal G_i]
 =\mathbb E[b_i(\theta_t)\mid\mathcal G_i].
\]
Here $Q_{ii}=q+\gamma$, $Q_{ij}=\gamma$ for $j\ne i$, and $b_i(\theta)=\theta_i$.
Let $\mu_j=\mathbb E[u_j]$. For $j\ne i$, the policy $u_j$ is a function of $\theta_{j,t}$,
which is independent of $\theta_{i,t}$, so
\[
\mathbb E[u_j\mid\mathcal G^{\mathrm{loc}}_{i,t}]
=\mathbb E[u_j]=\mu_j,
\]
while $\theta_{i,t}$ is
$\mathcal G^{\mathrm{loc}}_{i,t}$-measurable. The conditional normal equation
\eqref{eq:block_conditional_normal_equations} therefore gives
\[
(q+\gamma)u_i
+\gamma\sum_{j\ne i}\mu_j
-\theta_{i,t}=0.
\]
Taking expectations yields
\[
q\mu_i+\gamma\sum_j\mu_j=0.
\]
Summing over $i$ implies $\sum_j\mu_j=0$, hence every $\mu_i=0$ and
$u_i=\theta_{i,t}/(q+\gamma)$, which proves
\eqref{eq:optimal_local_shared}.

\emph{No-information regret.} Here $B=I$ and $\Sigma_S=\sigma^2I$, so by \eqref{eq:Rinf_general}, $R_\infty=\tfrac{\sigma^2}2\operatorname{tr}(Q^{-1})$. From the Sherman--Morrison formula above,

\[
\operatorname{tr}(Q^{-1})=\frac1q\Big(n-\frac{\gamma n}{q+\gamma n}\Big)=\frac{n\,(q+\gamma(n-1))}{q\,(q+\gamma n)},
\]

which gives \eqref{eq:aggregate_Rinf}.

\emph{Local regret.} By the Pythagorean form of Lemma~\ref{lem:radner_projection} (the orthogonality identity in Appendix~\ref{app:technical-projection}, with both $x^\star$ and $u^{\rm loc}$ mean zero), $2R_{\rm loc}=\|x^\star\|_Q^2-\|u^{\rm loc}\|_Q^2=2R_\infty-\|u^{\rm loc}\|_Q^2$. Since $u^{\rm loc}=\theta/(q+\gamma)$ and $\operatorname{tr}(Q)=n(q+\gamma)$,

\[
\|u^{\rm loc}\|_Q^2=\frac{\sigma^2\operatorname{tr}(Q)}{(q+\gamma)^2}=\frac{n\sigma^2}{q+\gamma}.
\]

Therefore

\[
R_{\rm loc}=\frac{n\sigma^2}{2}\left[\frac{q+\gamma(n-1)}{q(q+\gamma n)}-\frac{1}{q+\gamma}\right]=\frac{n\sigma^2}{2}\cdot\frac{(q+\gamma(n-1))(q+\gamma)-q(q+\gamma n)}{q(q+\gamma)(q+\gamma n)} .
\]

The numerator reduces to $\gamma^2(n-1)$, which gives \eqref{eq:aggregate_Rloc}.

\emph{Share.} Dividing \eqref{eq:aggregate_Rloc} by \eqref{eq:aggregate_Rinf}, the factors $n\sigma^2/2$ and $q(q+\gamma n)$ cancel, leaving \eqref{eq:aggregate_eta}.

\end{proof}

\section{Canonical spatial omission calculation}
\label{app:technical-canonical}
We expand the proof of Lemma~\ref{prop:canonical_eta_S} above,
including the equal-inverse-length case in the variance integral.
We use a continuous version of the field; observations on an interval can
equivalently be generated by a countable dense subset (including the center
at radius zero).
\begin{proof}
By spatial stationarity, take $z=0$ and suppress the time argument. Write $\beta_c\triangleq1/\ell_c$ and $\beta_s\triangleq1/\ell_s$ for the inverse lengths, so $w(u)=\tfrac{\beta_c}{2}e^{-\beta_c|u|}$ and $\operatorname{Cov}(\theta(u),\theta(u'))=\sigma^2e^{-\beta_s|u-u'|}$.

\textbf{Step 0 (omission is a variance ratio).} With $Q=qI$ the optimal policy given $\mathcal O_r$ is $\mathbb E[x^\star(0)\mid\mathcal O_r]$. Its expected loss is $\tfrac q2\,\mathbb E\big[\operatorname{Var}(x^\star(0)\mid\mathcal O_r)\big]$. For a Gaussian field the conditional variance is a deterministic number: it does not depend on the observed values (the Gaussian conditional covariance is independent of the observed values). Since $R_\infty=\tfrac q2\operatorname{Var}(x^\star(0))$,

\[
\eta_S(r)=\frac{\operatorname{Var}(x^\star(0)\mid\mathcal O_r)}{\operatorname{Var}(x^\star(0))}.
\]

\textbf{Step 1 (split into observed middle and two tails).} Write $x^\star(0)=M_r+T_r^++T_r^-$, where

\[
M_r=\int_{-r}^{r}w(u)\theta(u)\,du,\qquad T_r^+=\int_r^\infty w(u)\theta(u)\,du,\qquad T_r^-=\int_{-\infty}^{-r}w(u)\theta(u)\,du .
\]

The middle $M_r$ is computed from observed values, so it contributes no conditional variance.

\textbf{Step 2 (rewrite the right tail).} Substituting $u=r+s$,

\[
T_r^+=\frac{\beta_c}{2}e^{-\beta_c r}\,Y_r,\qquad Y_r\triangleq\int_0^\infty e^{-\beta_c s}\,\theta(r+s)\,ds .
\]

\textbf{Step 3 (spatial Markov property).} By the exponential-covariance model in Section~\ref{subsec:canonical_line_model}, the field is Markov in $z$. Given the observed segment $[-r,r]$, the right tail $\{\theta(u):u>r\}$ and the left tail $\{\theta(u):u<-r\}$ are conditionally independent, and each depends on the segment only through its adjacent endpoint $\theta(r)$ or $\theta(-r)$. Therefore

\[
\operatorname{Var}(T_r^+\mid\mathcal O_r)=\operatorname{Var}(T_r^+\mid\theta(r)),\qquad \operatorname{Var}(T_r^-\mid\mathcal O_r)=\operatorname{Var}(T_r^-\mid\theta(-r)),\qquad \operatorname{Cov}(T_r^+,T_r^-\mid\mathcal O_r)=0 .
\]

\textbf{Step 4 (the tail variance does not depend on $r$).} Define $C\triangleq\operatorname{Var}(Y_0\mid\theta(0))$. Shifting the field by $r$ does not change its law, so $\operatorname{Var}(Y_r\mid\theta(r))=C$ for every $r$. By reflection symmetry (the covariance depends only on $|u-u'|$), the left tail gives the same value. Hence

\[
\operatorname{Var}(x^\star(0)\mid\mathcal O_r)=2\Big(\frac{\beta_c}{2}\Big)^2e^{-2\beta_c r}\,C=\frac{\beta_c^2}{2}e^{-2\beta_c r}\,C .
\]

At $r=0$ the same argument applies with $\mathcal O_0=\sigma(\theta(0))$ and $M_0=0$, giving $\operatorname{Var}(x^\star(0)\mid\theta(0))=\tfrac{\beta_c^2}{2}C$. Dividing, $\eta_S(r)=\eta_S(0)\,e^{-2\beta_c r}=\eta_S(0)\,e^{-2r/\ell_c}$. The unknown constant $C$ cancels, so it never has to be computed.

\textbf{Step 5 (the constant $\eta_S(0)$).} Conditioning on the single Gaussian variable $\theta(0)$ gives $\operatorname{Var}(x^\star(0)\mid\theta(0))=\operatorname{Var}(x^\star(0))\big(1-\operatorname{Corr}(x^\star(0),\theta(0))^2\big)$ (the Gaussian conditional covariance is independent of the observed values), so $\eta_S(0)=1-\operatorname{Corr}^2$. Two moments are needed.

\emph{Covariance with the local value:}

\[
\operatorname{Cov}(x^\star(0),\theta(0))=\int_{\mathbb R}w(u)\,\sigma^2e^{-\beta_s|u|}\,du=\frac{\sigma^2\beta_c}{2}\int_{\mathbb R}e^{-(\beta_c+\beta_s)|u|}\,du=\frac{\sigma^2\beta_c}{\beta_c+\beta_s}.
\]

\emph{Variance of the target:} $\operatorname{Var}(x^\star(0))=\frac{\sigma^2\beta_c^2}{4}\,I$, where $I\triangleq\iint e^{-\beta_c(|u|+|v|)-\beta_s|u-v|}\,du\,dv$.

To evaluate $I$, define the inner integral $h(u)\triangleq\int_{\mathbb R}e^{-\beta_c|v|}e^{-\beta_s|u-v|}\,dv$. Then $h$ is even, so $I=2\int_0^\infty e^{-\beta_c u}h(u)\,du$. For $u\ge0$, split the $v$-integral at $0$ and $u$:

\[
h(u)=\underbrace{\frac{e^{-\beta_s u}}{\beta_c+\beta_s}}_{v<0}+\underbrace{\frac{e^{-\beta_c u}-e^{-\beta_s u}}{\beta_s-\beta_c}}_{0\le v\le u}+\underbrace{\frac{e^{-\beta_c u}}{\beta_c+\beta_s}}_{v>u}.
\]

(The middle term is read as $u\,e^{-\beta_c u}$ when $\beta_s=\beta_c$.) Multiply by $e^{-\beta_c u}$ and integrate over $u\ge0$ term by term:

\[
\frac{1}{(\beta_c+\beta_s)^2},\qquad \frac{1}{\beta_s-\beta_c}\Big(\frac{1}{2\beta_c}-\frac{1}{\beta_c+\beta_s}\Big)=\frac{1}{2\beta_c(\beta_c+\beta_s)},\qquad \frac{1}{2\beta_c(\beta_c+\beta_s)} .
\]

Summing and doubling,

\begin{equation}
I=2\left[\frac{1}{(\beta_c+\beta_s)^2}+\frac{1}{\beta_c(\beta_c+\beta_s)}\right]=\frac{2(2\beta_c+\beta_s)}{\beta_c(\beta_c+\beta_s)^2},\qquad \operatorname{Var}(x^\star(0))=\frac{\sigma^2\beta_c(2\beta_c+\beta_s)}{2(\beta_c+\beta_s)^2}.
\end{equation}

\emph{Combine.} With $\operatorname{Var}(\theta(0))=\sigma^2$,

\[
\operatorname{Corr}^2=\frac{\big(\sigma^2\beta_c/(\beta_c+\beta_s)\big)^2}{\sigma^2\cdot\sigma^2\beta_c(2\beta_c+\beta_s)/\big(2(\beta_c+\beta_s)^2\big)}=\frac{2\beta_c}{2\beta_c+\beta_s},\qquad \eta_S(0)=\frac{\beta_s}{2\beta_c+\beta_s}=\frac{\ell_c}{\ell_c+2\ell_s}.
\]
\end{proof}

\paragraph{Monotonicity in the decision-relevance length.}
For the positive-radius regime $vT>2\ell_s$, Theorem~\ref{thm:canonical_optimal_radius}
can be written as
\begin{equation}
r^\star=\frac12\int_{2\ell_s}^{vT}\frac{\ell_c}{\ell_c+s}\,ds.
\label{eq:rstar_integral_representation}
\end{equation}
With $vT$ and $\ell_s$ fixed, the integration limits do not change and
$\ell_c/(\ell_c+s)=1-s/(\ell_c+s)$ is strictly increasing in $\ell_c$
for every $s>0$. Thus $r^\star$ increases strictly with $\ell_c$ in this
regime, proving the fourth comparative static.

\section{Numerical reproducibility details}
\label{app:technical-numerics}
This section gives the settings and estimators for the four figures in
Section~\ref{sec:numerical-evaluation}. The crossover uses the exact coupled
fresh-local rule; the radius curves use scalar variance composition; and
the ring and four-geometry comparison use agentwise Gaussian conditioning with $Q=I$.
Normalized scalar and graph losses use analytical no-information regret
$R_\infty$. The crossover instead uses paired sample means.
Pointwise Monte Carlo means have 95\% normal intervals
$\bar L\pm1.96s_L/\sqrt N$, with sample standard deviation $s_L$.
Conditioning uses linear solves rather than explicit matrix inverses.

\begin{table}[H]
\centering
\caption{Settings for the four numerical illustrations. The shared-resource
model fixes $q=1$; graph covariances have average marginal variance one.}
\label{tab:numerical-parameters}
\small
\begin{tabular}{llll}
\toprule
Figure & principal sweep & fixed setting/size & samples \\
\midrule
\ref{fig:local-global-phase} & 9 $\gamma/q$ values in $[0,10]$ & $n=5,10,25,100$ & $30{,}000$/setting \\
\ref{fig:radius-regret} & $r/\ell_c\in[0,2]$ (1,001 points) & 3 regimes; $D/\ell_c=2$ & $50{,}000$/point \\
\ref{fig:decision-relevance} & $\ell_c=0.7,2,6$ & 64-node ring; $T=6$ & $12{,}000$/check \\
\ref{fig:general-graphs} & all integer radii; $T,\ell_c$ & four 64-node graphs & exact \\
\bottomrule
\end{tabular}
\end{table}

\paragraph{Fresh-local/stale-global crossover (Fig.~\ref{fig:local-global-phase}).}
We fix $q=1$, use $n\in\{5,10,25,100\}$, and sample independent
standard-normal state components ($\sigma^2=1$).
The nine coupling ratios are $\gamma/q\in\{0,0.05,0.1,0.2,0.5,1,2,5,10\}$.
At each setting, 30,000 independent current-state draws give paired no-information and fresh-local
losses, using $u_i=\theta_i/(q+\gamma)$.
Writing their means as $o=\widehat R_\infty$ and
$l=\widehat R_{\rm loc}$, the estimator is
$\widehat\rho^\star=\tfrac12\log[o/(o-l)]$.
Its 95\% delta-method interval is
$\widehat\rho^\star\pm1.96\sqrt{g^\top Cg/N}$, where $C$ is the sample
covariance of the paired losses and
\[
g=\left(-\frac{l}{2o(o-l)},\,\frac{1}{2(o-l)}\right)^\top.
\]
The paired covariance is retained; $o-l\le0$ is rejected rather than clipped.
The zero-coupling boundary is zero and is omitted from the logarithmic plot.

\paragraph{Three radius regimes (Fig.~\ref{fig:radius-regret}).}
The scalar calculation uses $\eta_S(r)=\eta_0e^{-2r/\ell_c}$,
$\tau(r)=r/v$, $\ell_c=1$, and $D=2$, with
$(\eta_0,vT)=(0.25,0.2),(0.7,4),(0.9,100)$.
The radius grid has spacing 0.002. At each point, 50,000 pairs of independent
standard normals are scaled by the square roots of the normalized delay and limited-scope terms in
\eqref{eq:R_decomposition}, then added
and squared. This checks scalar variance composition; it does not simulate
a spatial field.

\paragraph{Decision relevance on a ring (Fig.~\ref{fig:decision-relevance}).}
The graph and covariance are fixed on a 64-node ring, with $T=6$ and latency
0.12 per hop. For $\ell_c\in\{0.7,2,6\}$ the operator is
$K_{ij}=e^{-d_G(i,j)/\ell_c}/\|e^{-d_G(i,\cdot)/\ell_c}\|_2$.
With $Q=I$, exact Gaussian conditioning gives spatial omission; 12,000-sample
checks evaluate selected radii. Each operator has its own analytical
$R_\infty=\tfrac12\operatorname{tr}(K\Sigma_S K^\top)$.
Raw state correlation is computed separately from the fixed covariance.

The covariance in Eq.~\eqref{eq:graph_spectral_covariance} uses
$(\kappa_s,\nu)=(0.75,1.5)$ and is scaled to average marginal variance one.
Every integer radius from 0 through the ring diameter 32 is evaluated.
For each operator, Monte Carlo checks use radii
$\{0,r^\star,\min(r^\star+3,32),32\}$. Independent current-state Gaussian
draws are passed through the exact conditional policies; summed squared
prediction errors are divided by $\operatorname{tr}(K\Sigma_SK^\top)$
to estimate the omission share.

\paragraph{Finite noncanonical graphs (Fig.~\ref{fig:general-graphs}).}
Experiment 6 uses a 64-node path and cycle, an $8\times8$
grid without wraparound, and one connected random geometric graph.
For the latter, NumPy's \texttt{default\_rng(20260839)} draws 64 independent
uniform positions in $[0,1]^2$. Edges join pairs within Euclidean distance
0.22; the construction increases this threshold by 0.02 until connected
(0.22 suffices for this seed). Observation radii use hop distance in every
graph. Their diameters are 63, 32, 14, and 8, respectively.
There is no graph-ensemble average or Monte Carlo estimation in this experiment.

The covariance is proportional to $(0.8^2I+L_G)^{-1.5}$, scaled to average
marginal variance one. The target rows are
$K_{ij}=e^{-d_G(i,j)/\ell_c}/\|e^{-d_G(i,\cdot)/\ell_c}\|_2$ and $Q=I$.
For neighborhood $O_i=\mathcal N_r(i)$, the residual variance is
$K_i\Sigma_S K_i^\top-c_i^\top\Sigma_{S,O_iO_i}^{-1}c_i$,
where $c_i=\Sigma_{S,O_i,:}K_i^\top$; linear solves evaluate this expression.
The sum divided by $\operatorname{tr}(K\Sigma_S K^\top)$ gives $\eta_S(r)$.
Writing $\delta$ for the per-hop delay, so $\tau(r)=\delta r$,
we evaluate $1-e^{-2\delta r/T}(1-\eta_S(r))$ at every integer radius from zero
through the diameter, using reference $(T,\ell_c,\delta)=(5,2,0.12)$.

The plotted sweeps vary $T\in\{1.5,2.5,4,6,9,14\}$ or
$\ell_c\in\{0.7,1.2,2,3.5,6\}$ alone. The archive also retains
per-hop delays $\delta\in\{0.05,0.08,0.12,0.18,0.28\}$ and covariance parameters
$\kappa_s\in\{1.5,1,0.7,0.45,0.3\}$, each varied alone.
Configuration \path{numerics/config/experiment_06.yaml} records experiment
seed 20260828 and graph seed 20260839. The driver is
\path{numerics/experiments/exp06_general_graphs.py}, which calls
\texttt{exp06()} in \path{numerics/experiments/run_all.py}.
From the repository root, run \texttt{python scripts/verify\_graph\_restoration.py}
to rerun that function in scratch space and compare both processed CSVs and
all four raw graph archives under \path{numerics/results/}.
The audit checks positive-definite covariances, conditioning through an
independent residual-error calculation, and exhaustive integer-radius minima.
Its report is \path{docs/graph-restoration-verification.json}.

\paragraph{Reproduction of Figures~\ref{fig:local-global-phase}--\ref{fig:decision-relevance}.}
The \href{https://github.com/ANRGUSC/TINA}{public TINA repository} retains
the original experiment identifiers. Figures~\ref{fig:local-global-phase},
\ref{fig:radius-regret}, and \ref{fig:decision-relevance} correspond to
\path{numerics/config/experiment_02.yaml},
\path{numerics/config/experiment_03.yaml}, and
\path{numerics/config/experiment_05.yaml}, with seeds 20260824, 20260825,
and 20260827, respectively. Their drivers are
\path{numerics/experiments/exp02_local_global_crossover.py},
\path{numerics/experiments/exp03_radius_regret.py}, and
\path{numerics/experiments/exp05_decision_relevance.py}.
Run metadata, processed results, and selected raw arrays are stored under
\path{numerics/results/}; not every random draw is archived.
\path{docs/numerical_method.md} and \path{docs/experiment_parameters.md}
specify conditioning safeguards and environment versions;
\path{numerics/README.md} gives reproduction commands.


